\section{Approximability}\label{section4}
\subsection{Transformer can approximate finite steps of Turing-Machine}\label{section:approximation}

In this section, we present our theory showing that auto-regressive Transformers with a limited context window are universal approximators of any Turing computable functions. Although the simulation of a Turing machine does not guarantee that the model %\jx{Q: should we clarify that auto-regressive Transformers is equivalent to the language model in this paper??}
will always find the correct answer, it is a necessary condition for the language model to be able to effectively generate the expected text. %\jx{Q: should we say ``any target''? Also, did we define what target means, i.e., any given expected texts to be generated by LLM?}
By demonstrating that the Transformer's computational power is capable of approximating a Turing machine, we show that it reaches Turing Machine's limits %its theoretical limits \jx{Q:  is this the Turing machine's theoretical limits?} 
when handling complex sequential tasks.

We begin by encoding the computation of a corresponding Turing machine \( M \) that computes \( f \). Each computation step of \( M \) is represented as a configuration \( C_t \), encapsulating the current state, tape contents, and head position. These configurations form a sequence \( C_0, C_1, \dots, C_T \), where \( C_0 \) is the initial configuration based on the input \( x \), and \( C_T \), where: \( C_0 = {\color{black}\{x,q_0,\#\}}, \quad C_{t+1} = \delta(C_t), \) and  \(C_T = {\color{black}\{\cdot,q_{\text{accept}},\cdot\} }\) encodes the halting configuration producing \( f(x) \),  represents the halting configuration yielding \( f(x) \), {\color{black} where \(\#\) represents empty tape}.
The transition function \( \delta \) of \( M \) updates only a finite region of the tape based on the current state and symbol under the head:
  \(
  \delta: \Gamma^* {\color{black}\mathbb{Q}} \Gamma^* \rightarrow \Gamma^* {\color{black}\mathbb{Q}} \Gamma^*,
  \) {\color{black} where \(\Gamma*\) is empirical tape symbol space, and \(\mathbb{Q}\) is TM state space}


We first simulate the finite steps of a Turing Machine. The following lemmas show that several steps of the Turing machine can be simulated by an auto-regressive generative Transformer.

%\begin{lemma}\label{lemma:Turing-Transformer}
%For any deterministic Turing Machine \(\mathcal{M}\) that operates for \(S\) steps, there exists a Transformer model \(\mathcal{T}\) with finite parameters—specifically, finite number of layers \(L\), finite dimension per layer \(d\), finite acceptance window \(k\), and finite quantization levels \(Q\)—such that \(\mathcal{T}\) can asymptotically approximate the computation of \(\mathcal{M}\) for \(S\) steps. Furthermore, the maximum number of steps \(S_{\text{max}}\) that \(\mathcal{T}\) can approximate scales in \(\Theta(L\cdot d \cdot k \cdot \log(Q))\).
%\end{lemma}

\begin{lemma}\label{lemma:Turing-Transformer-Epsilon-Delta}
    Let \(\mathcal{M}\) be any deterministic Turing Machine that operates in \(S\) steps. For any \(\epsilon > 0\), there exists a Transformer model \(\mathcal{T}\) characterized by a finite number of layers \(L\), layer dimension \(d\), attention window size \(k\), and quantization levels \(Q\), such that for all computational steps \(s \leq S\), the state of \(\mathcal{T}\) approximates the state of \(\mathcal{M}\) at step \(s\) within an error bound \(\epsilon\). 
   {\color{black}\[
\forall \epsilon > 0, \, \exists \mathcal{T} \text{ with parameters } (L, d, k, Q) \text{ satisfying } 
\begin{cases}
 d \geq \log_2(|\mathbb{Q}|) + k \cdot \log_2(|\Gamma|), \\
Q \geq e^{\frac{C''' \cdot L \cdot d \cdot k}{\varepsilon}}
\end{cases}\]\(
\text{ such that } \forall s \leq S, \, d(H_s, \phi(C_s)) \leq \epsilon.
\)}
\end{lemma}

A key feature of the Turing Machine is its ability to store and access a large amount of information via its ``Tape''. When Transformers simulate Turing Machines through multi-round generation, they rely on the attention mechanism to store and retrieve information from previous generations. This suggests that Transformers can function as dynamic memory systems during the multi-round process, akin to the way a Turing Machine reads and writes on its tape.
\begin{lemma}\label{lemma:Turing-Transformer}
    
    The maximum number of computational steps \(S_{\text{max}}\) that \(\mathcal{T}\) can approximate while maintaining this error bound scales asymptotically as
    \[
        S_{\text{max}} \in \Theta\left(L \cdot d \cdot k \cdot \log(Q)\right).
    \]
\end{lemma}


Lemma \ref{lemma:Turing-Transformer} illustrates that a finite-precision, finite-depth, finite-width Transformer has only limited problem-solving capabilities. We demonstrate Lemma \ref{lemma:Turing-Transformer-Epsilon-Delta} and \ref{lemma:Turing-Transformer} in Appendix \ref{appendix:proof1} and \ref{appendix:b}. 

\subsection{Multi-Round Transformers are Turing-Machine Approximator}
We now consider the more far-reaching case: even if the Turing machine does not reach the halting condition within the Transformer's maximum generation window, we show that the Transformer can still approximate the simulation of the Turing machine through multiple rounds of generation. %\jx{Q: the previous sentence is problematic. Check again??}


\begin{theorem}[Approximability]\label{T:approximation}
    For any Turing-computable sequence-to-sequence function \( f: \Sigma^* \rightarrow \Sigma^* \), and for any error tolerance \( \epsilon > 0 \), there exists a multi-round sequence-to-sequence process utilizing an auto-regressive Transformer with a limited context window of size \( k \) that approximates \( f \) within error \( \epsilon \).
\end{theorem}

Theorem~\ref{T:approximation} can be inferred by induction directly from Lemma~\ref{lemma:Turing-Transformer-Epsilon-Delta}. The core idea for the proof is to consider that a Turing machine will halt within a sequence of finite length \(T\), by simulating the Turing Machine within \( R = \lceil T / s \rceil \) rounds and using induction for error propagation, the approximation error can be controlled within the tolerance. See detailed proof in Appendix~\ref{appendix:theorem43}.





%\jx{Once the error tolerance distribution is done, the induction seemed to be quite straightforward. Should we have some more interesting proof to show instead of this? For example, use the space to show the proof of other lemmas instead?} \chenhui{I am considering move this proof to appendix. I will add some discussion here.}


Thus, we conclude that any Turing-computable sequence-to-sequence function \( f: \Sigma^* \rightarrow \Sigma^* \) can be universally approximated by a multi-round sequence-to-sequence refinement process utilizing an auto-regressive Transformer with a limited context window of size \( k \), achieving the desired approximation within error \( \epsilon \).
Through multi-round generation, the Transformer moves beyond static one-shot input-output mappings and instead continuously adjusts its generation, similar to how human reasoning progresses. This dynamic computational ability is crucial for cognitive tasks as it allows the model to update its internal state and strategy during the process. This means that Transformers, beyond being powerful sequence generation models (e.g., for language translation), could potentially be applied to complex tasks such as cognitive reasoning, planning, and problem-solving.


%after \( R \) rounds, the final output \( y^{(R)} \) approximates \( f(x) \) within the total error tolerance \( \epsilon \) with the specified confidence \( 1 - \delta \). The sample complexity required to train the ARLM scales polynomially with \( 1/\epsilon \), \( 1/\delta \), and \( k \), ensuring that the process remains feasible within the PAC-learning framework.

%To ensure that the approximation holds with high probability, we utilize the union bound across all \( R \) rounds.  Each round \( r \) has a confidence level \( 1 - \delta_r = 1 - \delta / R \).
%By the union bound, the probability that any round \( r \) fails to meet its error bound is at most:
 % \[
  %\Pr\left( \bigcup_{r=1}^R \text{Failure in Round } r \right) \leq \sum_{r=1}^R \Pr(\text{Failure in Round } r) \leq \sum_{r=1}^R \delta_r = \delta.
  %\]
  %Therefore, with probability at least \( 1 - \delta \), all \( R \) rounds succeed in maintaining their respective error bounds.



















%\subsection{Multi-Round Vocabulary Composition is Universal Approximation}

%\begin{equation}
%    \label{eq:H_v}
%\mathcal{H}_{C'} = \{\phi_{i_1}  ... \phi_{i_m} | \phi_{i_1},...,\phi_{i_m}  \in C', m \in \mathbb{Z}_+\}.
%\end{equation}

%$C'$ here is the set of one-round language models and $\mathcal{H}_{C'}$ is the hypothesis function space of a sequence of language generation model on $C'$.











%\begin{theorem}\label{th:main-OP}



%Let $\Omega \subset \mathbb{R}^d$ be a compact domain. Then, there is a finite set $C' \subset \overline{\text{Diff}_0(\mathbb{R}^d)}$ such that the hypothesis space $\mathcal{H}_C'$ in Eq.~\eqref{eq:H_v} can be $C$-Universial Approximation for $\text{Diff}_0(\mathbb{R}^d)$.
%\end{theorem}

%\textcolor{red} {A constructive proof and more Lemma.}


